\chapter {Hamiltonian mechanics with dissipation}[Metriplectic dynamics]
\label{ch:mtpld}
Metriplectic dynamics is a geometric formulation of classical dynamics that
combines two fundamentally different kinds of evolution:
\begin{enumerate}
   \item Hamiltonian evolution, generated by an antisymmetric Poisson bracket,
       which describes reversible, conservative motion.
   \item Dissipative evolution, generated by a symmetric, positive semidefinite
       bracket, which describes irreversible relaxation and entropy production.
\end{enumerate}
The word metriplectic combines metric and symplectic (more precisely, Poisson) structures.

The central idea is that a physical system may undergo irreversible evolution
while still conserving its total energy. This is particularly useful for
describing viscous fluids, plasma collisions, friction, relaxation processes,
and nonequilibrium thermodynamics. The framework was developed notably by 
Morrison \cite{Morrison1986} in the 1980s. 

\section{The fundamental equations}
Let the state of a physical system be described by variables
\[
z=(z^1,\ldots,z^N).
\]
Ordinary Hamiltonian dynamics has the form
\[
\dot z^a=J^{ab}(z)\frac{\partial H}{\partial z^b},
\]
where \(H(z)\) is the Hamiltonian and \(J^{ab}\) is an antisymmetric Poisson
tensor:
\[
J^{ab}=-J^{ba}.
\]
Metriplectic dynamics introduces a second generator, entropy \(S(z)\), and a symmetric tensor \(M^{ab}(z)\):
\[
\boxed{
\dot z^a
=
J^{ab}\frac{\partial H}{\partial z^b}
+
M^{ab}\frac{\partial S}{\partial z^b}.
}
\]
The two tensors satisfy different algebraic properties:
\[
\begin{aligned}
J^{ab}&=-J^{ba},\\
M^{ab}&=M^{ba},\\
v_aM^{ab}v_b&\geq0.
\end{aligned}
\]
The first tensor defines the Poisson bracket
\[
\{F,G\}
=
\frac{\partial F}{\partial z^a}
J^{ab}
\frac{\partial G}{\partial z^b},
\]
while the second defines the metric bracket
\[
(F,G)
=
\frac{\partial F}{\partial z^a}
M^{ab}
\frac{\partial G}{\partial z^b}.
\]
Consequently, the evolution of any observable \(F(z)\) is
\[
\boxed{
\frac{dF}{dt}=\{F,H\}+(F,S).
}
\]
Here, \(J\) must satisfy the Jacobi identity, whereas \(M\) is symmetric and
positive semidefinite but need not define a nondegenerate Riemannian metric.

\subsection{ Why the construction satisfies thermodynamics}
The crucial additional requirements are the degeneracy conditions
\[
\boxed{
J^{ab}\partial_b S=0,
\qquad
M^{ab}\partial_b H=0.
}
\]
Equivalently,
\[
\{F,S\}=0,\qquad (F,H)=0
\]
for every observable \(F\).
These conditions have important physical consequences.
Conservation of total energy
Taking \(F=H\),
\[
\frac{dH}{dt}
=
\underbrace{\{H,H\}}_{0}
+
\underbrace{(H,S)}_{0}
=0.
\]
Thus,
\[
\boxed{\dot H=0.}
\]
The first term vanishes by antisymmetry, the second by degeneracy of the
dissipative bracket.
Entropy production
Taking \(F=S\),
\[
\frac{dS}{dt}
=
\underbrace{\{S,H\}}_{0}
+
(S,S).
\]
Since \(M\) is positive semidefinite,
\[
\boxed{
\dot S
=
\partial_a S\,M^{ab}\partial_b S
\geq0.
}
\]
Thus the first and second laws of thermodynamics are incorporated directly into
the equations of motion.

Notice an important distinction: the dissipative contribution does not destroy
energy; it redistributes energy while producing entropy.
For an isolated system, this is physically more appropriate than simply adding
a friction force to Hamilton's equations and allowing the total energy to
disappear.

\section{Example: a damped harmonic oscillator}
Consider the ordinary damped oscillator
\[
m\ddot q+\gamma\dot q+kq=0.
\]
Its mechanical energy is
\[
E_{\mathrm{mech}}
=\frac{p^2}{2m}+\frac{kq^2}{2},
\]
and
\[
\frac{dE_{\mathrm{mech}}}{dt}
=-\gamma\dot q^2\leq0.
\]
In ordinary Hamiltonian mechanics, this is problematic because mechanical
energy is not conserved.
But physically, the missing mechanical energy becomes heat.

We can account for this within metriplectic dynamics.
Step 1. Enlarge the state space
Introduce an entropy variable \(s\), describing an internal thermal reservoir:
\[
z=(q,p,s).
\]
Define the total energy
\[
\boxed{
H(q,p,s)
=
\frac{p^2}{2m}
+\frac{kq^2}{2}
+U(s),
}
\]
where \(U(s)\) is the internal energy of the reservoir, with temperature
\[
T(s)=\frac{dU}{ds}>0.
\]
Take the entropy generator to be \(S=s\).

Step 2. Introduce the two geometric structures
The Poisson tensor is
\[
J=
\begin{pmatrix}
0&1&0\\
-1&0&0\\
0&0&0
\end{pmatrix}.
\]
Thus \(q,p\) form a canonical pair, while \(s\) is a Casimir of the Poisson
structure.

For the dissipative part, choose
\[
\boxed{
M=\frac{\gamma}{T}
\begin{pmatrix}
0&0&0\\
0&T^2&-Tv\\
0&-Tv&v^2
\end{pmatrix},
\qquad v=\frac pm.
}
\]
This matrix is symmetric and positive semidefinite. Indeed, it is an outer product:
\[
M=\frac{\gamma}{T}
\begin{pmatrix}0\\T\\-v\end{pmatrix}
\begin{pmatrix}0&T&-v\end{pmatrix}.
\]
Furthermore,
\[
M\nabla H=0
\]
because
\[
\nabla H=(kq,v,T)
\]
and \(Tv-vT=0\).

Step 3. Obtain the evolution equations
Substituting into
\[
\dot z=J\nabla H+M\nabla S
\]
gives
\[
\boxed{
\begin{aligned}
\dot q&=\frac pm,\\
\dot p&=-kq-\gamma\frac pm,\\
\dot s&=\frac{\gamma}{T}\frac{p^2}{m^2}.
\end{aligned}
}
\]
The first two equations describe the familiar damped harmonic oscillator.
The third describes the production of entropy.

The mechanical energy decreases according to
\[
\dot E_{\mathrm{mech}}=-\gamma v^2,
\]
while the internal energy increases:
\[
\dot U=T\dot s=\gamma v^2.
\]
Hence
\[
\boxed{
\dot H=\dot E_{\mathrm{mech}}+\dot U=0,
\qquad
\dot S=\frac{\gamma v^2}{T}\geq0.
}
\]
We have reproduced the ordinary friction equation while explicitly satisfying
both laws of thermodynamics.

The essential point: The damped oscillator alone is an open subsystem. The
metriplectic formulation describes an enlarged, energetically closed model
containing both the oscillator and its heat reservoir.
This illustrates the physical purpose of the formalism particularly well.

\section{The geometric meaning: symplectic versus metric}
This is perhaps the most interesting aspect of the construction.
For a Hamiltonian system,
\[
X_H=J\,dH.
\]
The Poisson tensor maps the differential of the energy into a vector field
tangent to the energy surfaces.
Because \(J\) is antisymmetric,
\[
dH(X_H)=0.
\]
The resulting flow describes reversible motion along constant-energy surfaces.

A gradient flow instead has the structure
\[
Y_S=M\,dS.
\]
When \(M\) is positive semidefinite,
\[
dS(Y_S)\geq0.
\]
Consequently, the flow moves in directions of increasing entropy.
But in metriplectic dynamics, \(M\,dH=0\), so even the entropy-increasing part
must remain tangent to constant-energy surfaces.

There is a subtle point here: \(M\) is generally not an ordinary positive
definite metric. It has a nontrivial kernel containing \(dH\).
That degeneracy is not an inconvenience. It is exactly what guarantees energy
conservation.

Another geometric consequence concerns the Poisson structure. Because
\[
J\,dS=0,
\]
entropy is a Casimir. The Hamiltonian flow therefore remains on the symplectic
leaves of the Poisson manifold, while the dissipative evolution can move
between those leaves.
This explains why metriplectic dynamics naturally uses noncanonical, degenerate Poisson structures.
In particular, a nonconstant entropy Casimir cannot exist on an ordinary
connected symplectic manifold with a nondegenerate Poisson tensor. One must
enlarge or reduce the description, or work with a genuinely degenerate Poisson
structure.

\section{Relation to other formulations of mechanics}
\begin{tabular}{|c|c|c|}
 \hline
    Framework&	Generators&	Basic feature\\ \hline
    Hamiltonian	&Energy \(H\)	&Reversible, energy conserving\\ \hline
    Gradient flow&	Potential \(S\) or \(V\)&	Monotonic relaxation\\\hline
    Metriplectic&	Energy \(H\), entropy \(S\)&	Energy conserved, entropy increases\\\hline
    Contact Hamiltonian&	Contact Hamiltonian&	Geometric description of some dissipative systems\\\hline
    GENERIC&	Energy \(E\), entropy \(S\)	&Reversible–irreversible thermodynamic coupling\\\hline
\end{tabular}

A particularly close relative is GENERIC (General Equation for Non-Equilibrium
Reversible–Irreversible Coupling), developed by Grmela and Öttinger
\cite{GrmelaÖttinger1997},\cite{ÖttingerGrmela1997}. In its
standard form, it has essentially the same two-generator structure and
degeneracy conditions, though its development emphasizes nonequilibrium
thermodynamic modeling and constitutive relations. 

There is also an important limitation: the metric tensor \(M\) is not uniquely
determined by energy conservation and entropy production. These requirements
constrain the dissipative dynamics but do not determine its magnitude or
detailed physical mechanism. Friction coefficients, transport coefficients, and
kinetic collision operators must still be supplied by additional physics.

An interesting newer development is the metriplectic four-bracket formulation
introduced by Morrison and Updike \cite{MorrisonUpdike2024}.
Instead of starting from an arbitrary binary symmetric bracket, it constructs
dissipative dynamics using a four-argument bracket with algebraic symmetries
analogous to those of the Riemann curvature tensor.
The purpose is to provide a more geometrically systematic construction of
admissible dissipation operators. This makes the relation between metriplectic
structures and differential geometry considerably deeper than the simple
combination of a symplectic form and a Riemannian metric. 

\paragraph{A question worth pursuing}
From the standpoint of classical mechanics, the
deepest issue is not merely how to append dissipation to Hamilton's equations,
but: Why should reversible mechanics be governed by an antisymmetric geometric
structure, while irreversible mechanics is governed by a symmetric one? In
Hamiltonian mechanics, the antisymmetric structure comes naturally from the
Legendre transformation and the symplectic geometry of phase space. In
metriplectic mechanics, however, the symmetric structure is not fixed by the
symplectic geometry or by the first and second laws alone. This raises a
substantive question: Can the dissipative metric bracket be derived from an
underlying microscopic Lagrangian or Hamiltonian theory by eliminating degrees
of freedom, rather than introduced phenomenologically? That question connects
metriplectic dynamics to statistical mechanics, coarse-graining, the
Mori–Zwanzig projection formalism, and the emergence of irreversible behavior
from reversible microscopic equations.
